\documentclass[11pt,aspectratio=169]{beamer}
%\documentclass[11pt,aspectratio=169,handout]{beamer}

\usepackage{amsmath}
\usepackage{amsfonts}
\usepackage{amssymb}
\usepackage{graphbox}
\usepackage{sgamevar}
\usepackage{pgfplots}
\usepackage{tikz}
\usepackage{pstricks,pst-node}
\usepackage{bigdelim}
\usepackage{qrcode}
\usepackage[absolute,overlay]{textpos}
\usepackage[ruled,vlined]{algorithm2e}



% alg
\SetKwFor{Repeat}{repeat}{}{}
\SetKwComment{Comment}{$\triangleright$ }{}
\SetCommentSty{textsf}
% pgfplot settings
\usepgfplotslibrary{fillbetween}

% tikz settings
\tikzset{
    invisible/.style={text opacity=1,opacity=0},
    visible on/.style={alt=#1{}{invisible}},
    alt/.code args={<#1>#2#3}{%
      \alt<#1>{\pgfkeysalso{#2}}{\pgfkeysalso{#3}}
    },
  }
\newcommand{\nebox}[2][]{\tikz[baseline=(h.base)]\node[rounded corners,rectangle,draw,line width=0.7pt,text depth=-4.9pt,#1] (h) {#2};}
\usetikzlibrary{arrows}
% Customized colors

\definecolor{ashgrey}{rgb}{0.7, 0.75, 0.71}
\definecolor{lightblue}{RGB}{140, 201, 247}
\definecolor{darkgreen}{RGB}{113, 160, 55}
\definecolor{darkblue}{RGB}{78, 103, 200}
\definecolor{darkpurple}{RGB}{112, 48, 160}

% Hyperlink settings

\hypersetup{colorlinks,urlcolor=darkgreen}

%Beamer settings

\setbeamertemplate{navigation symbols}{} 
\setbeamertemplate{footline}[frame number]
\setbeamertemplate{itemize items}[circle]
\setbeamertemplate{section in toc}[sections numbered]
\setbeamertemplate{subsection in toc}[subsection numbered]

\setbeamercolor{section in toc}{fg=blue}

\AtBeginSection[ ]{
 \begin{frame}{Outline}
  \hypersetup{linkcolor=black}
  \tableofcontents[currentsection]
 \end{frame}
}

% Graphics settings

\graphicspath{{../Figures/}}

% PGFPlot settings

\pgfplotsset{compat = newest}

% Gamesvar settings

\setlength{\arrayrulewidth}{0.91pt}
\renewcommand{\gamestretch}{1.68}
\def\stackedpayoffs#1#2{%
 \begin{array}{c}#1\\[2mm]#2\end{array}
}

% Math

\DeclareMathOperator*{\argmax}{argmax}
\DeclareMathOperator*{\argmin}{argmin}

\newcommand\given[1][]{\:#1\vert\:}
\newtheorem{proposition}{Proposition}

% New environments

\newenvironment{itemizes}[1][1em]{
 \vspace{#1}
 \begin{itemize}
 \setlength{\itemsep}{#1}
}{
 \end{itemize}
}


% Shared title frame

\title{Game-theoretic \\ Foundations of Multi-agent Systems}

\author{Seyed Majid Zahedi}
\titlegraphic{\vspace{-4.2em} \includegraphics[height=5.8em]{Logos/logo2}}

\date{} 
\subject{Lecture 1} 
\logo{\includegraphics[height=5.6em]{Logos/logo1}}
\subtitle{\vspace{2.1em}Lecture 6: Repeated Games}

\begin{document}

 \begin{frame}[plain]
  \titlepage
 \end{frame}
 
 
 \section{Finitely Repeated Games}
  \begin{frame}{Repeated Games}
   \begin{itemize}[<+->]
   \setlength{\itemsep}{1.2em}
    \item In a (typical) repeated game:
    \begin{itemize}[<.->]
     \item Agents play a given game (aka. \alert{stage game})
     \item Then, they get their utilities
     \item And, they play again $\dots$
    \end{itemize}
    \item Can be repeated \alert{finitely} or \alert{infinitely} many times
    \item Really, an extensive form game
    \begin{itemize}[<.->]
    \item Would like to find subgame-perfect equilibria
    \end{itemize}
   \end{itemize}
  \end{frame}
  
  
  \begin{frame}{Repeated Games (cont.)}
   \begin{itemize}[<+->]
   \setlength{\itemsep}{1.2em}
    \item One subgame-perfect equilibrium:
    \begin{itemize}[<.->]
     \item Keep repeating some Nash equilibrium of the stage game 
     \item \alert{Memoryless} strategy, called a \alert{stationary strategy}
    \end{itemize}
    \item But are there other equilibria?
    \begin{itemize}[<.->]
     \item Strategy space of repeated game is much richer than that of stage game
    \end{itemize}
   \end{itemize}
  \end{frame}


  \begin{frame}{Key Questions}
   \begin{itemize}
   \setlength{\itemsep}{1.2em}
    \item Do agents see what the other agents played earlier?
    \item Do they remember what they knew?
    \item Given utility of each stage game, what is the utility of the entire repeated game?
   \end{itemize}
  \end{frame}
  
  
  \begin{frame}{Finitely Repeated Games (with Perfect Monitoring)}
   \begin{itemize}[<+->]
   \setlength{\itemsep}{1.2em}
    \item Agents play stage game $G$ for $R$ rounds
    \item At each round, outcomes of all past rounds are observed by all agents
    \item Agents' overall utility is sum of \alert{discounted utilities} at each round
    \begin{itemize}[<.->]
     \item Discount factor is $ 0 \le \delta \le 1$
     \item Game is denoted by $G^{R}(\delta)$
    \end{itemize}
    \item Given sequence of utilities $u_i^{(1)}, ..., u_i^{(R)}$, $u_i = \sum_{r=1}^{R}\delta^{r-1} u_i^{(r)}$
   \end{itemize}
  \end{frame}
 
 
  \begin{frame}{Example: Finitely Repeated Prisoner's Dilemma}
   \begin{itemize}[<+->]
    \item Two agents play Prisoner's Dilemma for $R$ rounds ($\delta = 1$)
    \begin{center}
     \hspace{-3.5em}
     \begin{game}{2}{2}
      		\> D			\> C			\\
      D		\> $-2,-2$	\> $-4,-1$	\\
      C		\> $-1,-4$	\> $-3,-3$
     \end{game}
    \end{center}
    \vspace{0.7em}  
    \item Starting from last round, (C, C) is dominant strategy
    \item Hence, in second-to-last round, there is no way to influence what will happen
    \item So, (C, C) is dominant strategy at this round as well
    \item The unique SPE is (C, C) at each round
   \end{itemize}
  \end{frame}
  
  \begin{frame}{SPE in Finitely Repeated Games}
   \alert{[Theorem]}
   \begin{itemize}
    \item If stage game $G$ has unique strategy equilibrium $s^{*}$, then $G^{R}(\delta)$ has unique SPE in which $s^{(r)} =s^{*} $ for all $r = 1, ..., R$, regardless of history
   \end{itemize}
   \vspace{1em}   
   \alert{[Proof]}
   \begin{itemize}
    \item By backward induction, at round $R$, we have $s^{(R)} = s^{*}$
    \item Given this, then we have $s^{(R - 1)} = s^{*}$, and continuing inductively, $s^{(r)} = s^{*}$ for all $r = 1, ..., R$, regardless of history
   \end{itemize}
  \end{frame}

  \begin{frame}{SPE: Example I}
   \begin{itemize}[<+->]
    \item Two agents play the following game for 2 rounds ($\delta = 1$)
    \begin{center}
     \hspace{-3.5em}
     \begin{game}{3}{3}
      		\> D1		\> D2		\> C			\\
      D1		\> $4, 4$	\> $1, 1$	\> $6, 0$	\\
      D2		\> $1, 1$	\> $2, 2$	\> $6, 0$	\\
      C		\> $0, 6$	\> $0, 6$	\> $5, 5$
     \end{game}
    \end{center}
    \vspace{0.7em}  
    \item Consider the following strategy:
    \begin{itemize}[<.->]
     \item In round 1, cooperate;
     \item In round 2, if someone defected in round 1, play D2; otherwise, play D1
    \end{itemize}
    \item If both agents play this, is that SPE?
   \end{itemize}
  \end{frame}  
 

  \begin{frame}{SPE: Example II}
   \begin{itemize}[<+->]
    \item Two agents play the following game for 2 rounds ($\delta = 1$)
    \begin{center}
     \hspace{-3.5em}
     \begin{game}{3}{3}
      		\> D			\> Crazy		\> C			\\
      D		\> $4, 4$	\> $1, 0$	\> $6, 0$	\\
      Crazy	\> $0, 1$	\> $0, 0$	\> $0, 1$	\\
      C		\> $0, 6$	\> $1, 0$	\> $5, 5$
     \end{game}
    \end{center}
    \vspace{0.7em}  
    \item What are the subgame perfect equilibria?
    \item Consider the following strategy:
    \begin{itemize}[<.->]
     \item In round 1, cooperate;
     \item In round 2, if someone played D or Crazy in round 1, play Crazy; otherwise, play D
    \end{itemize}
    \item If both agents play this, is that NE (not SPE)?
   \end{itemize}
  \end{frame} 
  
  
  \begin{frame}{TSPE: Example III}
   \begin{itemize}
    \item  If $G$ has multiple equilibria, then $G^{R}(\delta)$ does not have unique SPE
    \item Consider following example
    \begin{center}
     \hspace{-3.5em}
     \begin{game}{3}{3}
      	\> x			\> y			\> z			\\
      x	\> $3, 3$	\> $0, 4$	\> $-2, 0$	\\
      y	\> $4, 0$	\> $1, 1$	\> $-2, 0$	\\
      z	\> $0, -2$	\> $0, -2$	\> $-1, -1$
     \end{game}
    \end{center}
    \vspace{0.7em}
    \item Stage game has two pure NE: (y, y) and (z, z)
    \item Consider the following policy:
    \begin{itemize}
     \item Play x in first round
     \item Play y in second round if opponent played x; otherwise, play z
    \end{itemize}
    \item Is both agents playing this SPE?
   \end{itemize}  
  \end{frame}

 \section{Infinitely Repeated Games}
 
  \begin{frame}{Utilities in Infinitely Repeated Games}
   \begin{itemizes}[2em]
    \item \alert{Limit-average utility}: 
    $$u_i = \lim_{R\rightarrow\infty}\frac{\sum_{r=1}^{R}u_i^{(r)}}{R}$$
    \item \alert{Future-discounted utility}: 
    $$u_i = (1 - \delta) \sum_{r=1}^{\infty}\delta^{r-1} u_i^{(r)},$$
    for some $0 \ge \delta < 1$
   \end{itemizes}
  \end{frame}
  
  
  \begin{frame}{Subgame Perfection in Infinitely Repeated Games}
   \begin{itemize}
   \setlength{\itemsep}{2em}
    \item \alert{One-shot deviation} from strategy $s$ means deviating from $s$ in single stage and conforming to it thereafter
    \item Strategy profile $s^*$ is SPE \alert{if and only if} there are no \alert{profitable} one-shot deviation for \alert{each subgame} and \alert{every agent}
    \item This follows from principle of optimality of \alert{dynamic programming}
    \item This applies to finitely repeated games as well
   \end{itemize}
  \end{frame}

  \begin{frame}{Trigger Strategies (TS)}
   \begin{itemize}
    \item Agents get \alert{punished} if they deviate from agreed profile
    \item In \alert{non-forgiving} TS (or grim TS), punishment continues forever
    $$ 
    s_i^{(t)}=\begin{cases}
                  s_i^* &\text{if } s^{(r)} = s^* \;\; \forall r<t,\\
                  \underline{s}_{i}^{j} &\text{otherwise}\\
                 \end{cases}
    $$
    \item Here, $s^{*}$ is agreed profile, and $\underline{s}_{i}^{j}$ is punishment strategy of $i$ against agent $j$
    \item Single deviation by $j$ triggers agent $i$ to switch to $\underline{s}_{i}^{j}$ \alert{forever}
   \end{itemize}
  \end{frame}


  \begin{frame}{Example: Infinitely Repeated Prisoner's Dilemma} 
   \begin{columns}
    \begin{column}{0.7\textwidth}
     \begin{itemize}[<+->]
     \setlength{\itemsep}{1.2em}
      \item Consider \alert{trigger} strategy:
      \begin{itemize}[<.->]
       \item Cooperate as long as everyone cooperates
       \item Once a player defects, defect \alert{forever}
      \end{itemize}
      \item Is both agents playing this SPE?
      \item Does it depend on $\delta$?
     \end{itemize}
    \end{column}
    \begin{column}{0.3\textwidth}
     \begin{center}
      \begin{game}{2}{2}
      		\> D			\> C			\\
       D		\> $-2,-2$	\> $-4,-1$	\\
       C		\> $-1,-4$	\> $-3,-3$
      \end{game}
     \end{center}
    \end{column}
   \end{columns}
  \end{frame}
  
  
  \begin{frame}{Trigger Strategy for Infinitely Repeated Prisoners' Dilemma}
   \begin{itemize}[<+->]
   \setlength{\itemsep}{0.7em}
    \item We can use one-stage deviation principle
    \item There are two types of subgames:
    \begin{itemize}[<.->]
     \item Type 1: Both agents cooperated so far
     \item Type 2: At least one agent defeated in the past
    \end{itemize}
    \item Type-1 subgames: (D is best response to D)
    \begin{itemize}
     \item Utility from no deviation: $(1 - \delta) (-2 - 2\delta - 2\delta^{2} + \dots) = - 2$
     \item Utility from on-shot deviation: $(1 - \delta) (-1 + (-3\delta -3\delta^{2} + \dots) = - (1 - \delta ) - 3\delta$
     \item Deviation is not beneficial if $\delta  \geq  1/2$
    \end{itemize}
    \item Type-2 subgames: (C is best response to C)
    \begin{itemize}
     \item Other agents will always play C, thus C is best response
    \end{itemize}  
   \end{itemize}
  \end{frame}    
    
  \begin{frame}{Tit-for-tat Strategy}
   \begin{itemize}[<+->]
   \setlength{\itemsep}{1.2em}
    \item Consider \alert{tit-for-tat} strategy:
    \begin{itemize}[<.->]
     \item Cooperate in 1st round
     \item Then, do whatever other agent did in previous round
    \end{itemize}
    \item Is both agents playing this NE?
    \item Is both agents playing this SPE?
    \item What about one playing TFT and other trigger?
   \end{itemize}
  \end{frame}
  

  \begin{frame}{Remarks}
   \begin{itemize}
   \setlength{\itemsep}{2em}
    \item If $s^{*}$ is NE of $G$, then ``each agent plays $s_i^{*}$'' is SPE of $G^{R}(\delta)$
    \begin{itemize}
     \item Future play of other agents is independent of how each agent plays
     \item Optimal play is to maximize current utility, i.e., play static best response
    \end{itemize} 
    \item Sets of equilibria for finite and infinite horizon versions can be \alert{different}
    \begin{itemize}
     \item Multiplicity of equilibria in repeated prisoner's dilemma only occurs at $R = \infty$
     \item For any finite $R$ (thus for $R \rightarrow \infty$), repeated prisoners' dilemma has unique SPE
    \end{itemize}
   \end{itemize}
  \end{frame}


  \begin{frame}{Repetition Could Lead to Bad Outcomes}
   \begin{itemize}
   \setlength{\itemsep}{0.5em}
    \item Consider the following game
    \begin{center}\scriptsize
     \hspace{-3.5em}
     \begin{game}{3}{3}
      	\> x 		\> y			\> z			\\
      x	\> $2, 2$	\> $2, 1$	\> $0, 0$	\\
      y	\> $1, 2$	\> $1, 1$	\> $-1, 0$	\\
      z	\> $0, 0$	\> $0, -1$	\> $-1, -1$
     \end{game}
    \end{center}
    \vspace{0.7em}
    \item Strategy x strictly dominates y and z for both agents
    \item Unique NE of stage game is (x, x)
    \item If $\delta \geq 1/2 $, this game has SPE in which (y, y) is played in every round
    \item It is supported by slightly more complicated strategy than grim trigger
    \begin{itemize}
     \item I. Play y in every round unless someone deviates, then go to II
     \item II. Play z. If no one deviates go to I. If someone deviates stay in II
    \end{itemize}
   \end{itemize}
  \end{frame}
  
  
 \section{Folk Theorem}  
  
  \begin{frame}{Characterizing NE of Infinitely Repeated Games}
   \begin{itemize}[<+->]
   \setlength{\itemsep}{0.7em}
    \item Characterizing all equilibrium strategy profiles might be challenging
    \item Instead, we can characterize utilities obtained in them
    \item Such utilities must be \alert{feasible}
    \begin{itemize}[<.->]
     \item There must be outcomes of game such that agents, on average, get these utilities
    \end{itemize}
    \item They must also be \alert{enforceable}
    \begin{itemize}[<.->]
     \item Deviation should lead to punishment that outweighs benefits of deviation
    \end{itemize}
    \item \alert{Folk theorem} states that utility vector can be realized by some NE iff it is both feasible and enforceable
   \end{itemize}
  \end{frame}
  
  \begin{frame}{Feasible Utilities: Formal Definition}
   \begin{itemize}[<+->]
   \setlength{\itemsep}{1.2em}
    \item \alert{Utility profile} $u = (u_1,u_2,\dots,u_n)$ is \alert{feasible} if there exist \alert{rational}, \alert{non-negative} values $\{\alpha_a\}$ such that for all $i$, $u_i = \sum_{a\in A} \alpha_a u_i(a)$, with $\sum_{a\in A} \alpha_a = 1$
    \item You could think of feasible utilities as \alert{convex hull} of possible outcomes:
    $$U = \text{Conv}\{ u \in  \mathbb{R}^{\mid N \mid} \mid \text{ there exists } a \in A \text{ such that } u(a) = u \}$$
    \item Note that $U \neq  \{ u \in  \mathbb{R} ^{\mid N \mid}  \mid \text{ there exists } s \in S \text{ such that } u(s) = u \}$
   \end{itemize}
  \end{frame}
  
  
  \begin{frame}{Feasibility: Example}
   \begin{center}
    \hspace{-2em}
    \begin{game}{2}{2}
     		\> Left		\> Right		\\
     Left	\> $2, 2$	\> $0, 3$	\\
     Right	\> $3, 0$	\> $1, 1$
    \end{game}
   \end{center}
   \vspace{1em}
   \begin{itemize}[<+->]
    \item Utility vector (2, 2) is feasible as it is one of outcomes of game
    \item Utility vector (1, 2.5) is feasible as agents can alternate between (2, 2) and (0, 3)
    \item What about (0.5, 2.75)?
    \item What about (3, 0.1)?
   \end{itemize}
  \end{frame}


  \begin{frame}{Enforceable and Individually Rational Utilities}
   \begin{itemize}[<+->]
   \setlength{\itemsep}{1.2em}
    \item Recall \alert{minmax value} of agent $i$: 
    $$\underline{v}_i = \underset{s_{-i}}{\min} \; \underset{s_{i}}{\max} \; u_{i}(s_i,s_{-i})$$
    \item Utility profile $u \in  \mathbb{R} ^{\mid N \mid}$  is \alert{individually rational} if $u_{i}  \ge \underline{v}_i$ for all $i$
    \item Utility profile $u = (u_1,u_2,\dots,u_n)$ is \alert{enforceable} if it is individually rational
   \end{itemize}
  \end{frame}
  
  
  \begin{frame}{Nash Folk Theorem}
   \begin{itemize}[<+->]
   \setlength{\itemsep}{1.2em}
    \item Consider infinitely repeated game $G$ played by agents with \alert{average utilities}
    \item If $u$ is utility profile for any NE of repeated $G$, then $u_i$ is enforceable for all $i$
    \item If $u$ is both feasible and enforceable, then $u$ is utility profile for some NE of $G$
    \item Folk theorem can be stated for agents with discounted utilities as well
   \end{itemize}
  \end{frame}
  
  
  \begin{frame}{Problems with Nash Folk Theorem}
   \begin{itemize}[<+->]
    \setlength{\itemsep}{1em}
    \item Any feasible and enforceable utility can be achieved (for \alert{patient enough} agents)
    \item Enforcement is often done by grim trigger strategy
    \begin{itemize}
     \item Play certain strategy as long as no one deviates
     \item If some agent $j$  deviates, then play minmax strategy against that agent thereafter
    \end{itemize}
    \item NE involves non-forgiving TS which may be costly for punishers
    \item NE may include \alert{non-credible threats}
    \item NE may not be subgame perfect
   \end{itemize}
  \end{frame}
  
  
  \begin{frame}{Example}
   \begin{center}
    \hspace{-0.2em}
    \begin{game}{2}{2}
     	\> L			\> R				\\
     U	\> $6, 6$	\> $ 0, -100$	\\
     D	\> $7, 1$	\> $ 0, -100$
    \end{game}
   \end{center}
   \vspace{1em}
   \begin{itemize}
    \item Unique NE in this game is (D, L)
    \item Minmax values are given by $\underline{v}_{1} = 0 $ and $\underline{v}_{2} = 1 $ 
    \item Minmax strategy against agent 1 requires agent 2 to play R
    \item R is strictly dominated by L for agent 2
   \end{itemize}
  \end{frame}
  
 \section{Repeated Games with Imperfect Monitoring}
 
  \begin{frame}{Motivation}
   \begin{itemizes}[1.5em]
    \item So far, we assumed that agents observe actions of others at each round of game
    \item Next, we consider games where agents' actions may not be directly observable
    \item We assume that agents observe only an \alert{imperfect signal} of stage game actions
   \end{itemizes}
  \end{frame}   
 
  
  \begin{frame}{Example: Cournot Competition with Noisy Demand\\
  \vspace{-0.2cm}{\scriptsize [Green and Porter, Non-cooperative Collusion under Imperfect Price Information, 1984]}}
   \begin{itemizes}
    \item Firms set production levels $q_1^{(r)}, \dots, q_n^{(r)}$ \alert{privately} at round $r$
    \item Firms do not observe each others' output levels
    \item Market demand is \alert{stochastic}
    \item Market price depends on total production and market demand
    \item Low price could be due to high production or low demand
    \item Firms utility depends on their own production and market price
   \end{itemizes}
  \end{frame}


  \begin{frame}{Model}
   \begin{itemizes}
    \item We focus on game with \alert{public information}
    \item At each round, \alert{all agents} observe some \alert{public outcome}
    \item Let $y^{(r)} \in Y$ denote publicly observed outcome at round $r$
    \item Each action profile $a$ induces \alert{probability distribution} over $y$
    \item Let $\pi (y, a)$ denote probability distribution of y under action profile $a$
    \item Public information at round $r$ is $h^{(r)}=(y^{(1)}, \dots, y^{(r-1)})$
    \item Strategy of agent $i$ is \alert{sequence of maps} $s_i^{(r)}: h^{(r)} \rightarrow S_i$ 
   \end{itemizes}
  \end{frame}
  
  \begin{frame}{Model (cont.)}
   \begin{itemize}[<+->]
    \item Agents utility depends \alert{only} on their own action and public outcome
    \item Dependence on actions of others is through their effect on distribution of $y$
    \item Agent $i$'s \alert{realized} utility at round $r$ is $u_i(a_i^{(r)}, y^{(r)})$
    \item Agent $i$'s expected stage utility is
    $$
    u_i(a) =\sum_{y\in Y}\pi(y,a)u_i(a_i,y) 
    $$
    \item Agent $i$'s average discounted utility when sequence $\{ a^{(t)}\}$ is played is
    $$
     (1 - \delta)\sum_{r=1}^{\infty}\delta^{r-1} u_i(a^{(r)}) 
    $$
   \end{itemize}
  \end{frame}
  
  
  \begin{frame}{Simpler Example: Noisy Prisoner's Dilemma}
   \begin{itemize}[<+->]
   \setlength{\itemsep}{0.5em}
    \item Prisoners do not observe each others actions, instead, they observe signal $y$
    \begin{itemize}
     \item $u_{1}(D, y) = 1 + y $   \hspace{1cm} $u_{1}(C, y) = 4 + y $ 
     \item $u_{2}(D, y) = 1 + y $   \hspace{1cm} $u_{2}(C, y) = 4 + y $ 
    \end{itemize}
    \item Signal $y$ is defined by cont. random variable $X$ with CDF $F(x)$ and $\mathbb{E}[X] = 0$
    \begin{itemize}
     \item If $a = (D, D)$, then $y = X$
     \item If $a = (D, C)$ or $(C, D)$, then $y = X - 2$
     \item If $a = (D, D)$, then $y = X - 4$ 
    \end{itemize}
    \item Normal-form stage game is
    \begin{center}\footnotesize
     \hspace{-6em}
     \begin{game}{2}{2}
      	\> D					\> C						\\
      D	\> $1 + X, 1 + X$	\> $ -1 + X, 2 + X$	\\
      C	\> $2 + X, -1 + X$	\> $ X, X$
     \end{game}
    \end{center}
   \end{itemize}
  \end{frame}

  
  \begin{frame} {Trigger-price Strategy}
   \begin{itemizes}[1.2em]
    \item Consider following trigger strategy
    \begin{itemize}
     \item (I) - Play $(D, D)$ until $y \le y^{*}$, then go to (II)
     \item (II) - Play$ (C, C)$ for $R$ rounds, then go back to (I)
    \end{itemize}
    \item Notice that strategy is stationary and symmetric
    \item Also notice that punishment uses \alert{NE of stage game}
    \item We can choose $y^{*}$ and $R$ such that this strategy profile is SPE
   \end{itemizes}
  \end{frame}


  \begin{frame} {Trigger-price Strategy (cont.)}
   \begin{itemizes}
    \item We use one-shot deviation principle
    \item Deviation in (II) is obviously not beneficial
    \item In (I), if agents do not deviate, their \alert{expected utility} is
    $$
    v = (1-\delta)\left((1+0)+\delta\left(F(y^*)\delta^{R} v + (1-F(y^*)) v \right)\right)
    $$
    \item From this, we obtain
    $$
    v = \frac{1-\delta}{1-\delta(1-\delta)\big(1-F(y^*)(1-\delta^{R})\big)}
    $$    
   \end{itemizes}
  \end{frame}

  \begin{frame}{Trigger-price Strategy (cont.)}
   
   \begin{itemize}
    \item If some agent deviates in (1), then her expected utility is
    $$
    v_d = (1-\delta)\left((2+0)+\delta\left(F(y^*+2)\delta^{R} v + (1-F(y^*+2)) v \right)\right)
    $$
    \item Deviation provides immediate utility, but increases probability of entering (II)
    \item To have SPE, we mush have $v \geq v_d$ which means
    $$
    v \ge \frac{2(1-\delta)}{1-\delta(1-\delta)\big(1-F(y^*+2)(1-\delta^{R})\big)}
    $$
    $$
    \Rightarrow F(y^*+2)-2F(y^*) \ge \frac{1-\delta(1-\delta)}{\delta(1-\delta)(1-\delta^{R})}
    $$
    \item  Any $R$ and $y^{*}$ that satisfy this constraint construct SPE
    \item \alert{Best trigger-price strategy} can be found by maximizing $v$ s.t. this constraint
   \end{itemize}
 
  \end{frame}  
  
  
  \begin{frame}{Acknowledgment}
   \begin{itemize}
   \setlength{\itemsep}{1em}
    \item This lecture is a slightly modified version of ones prepared by
    \begin{itemize}
     \item Asu Ozdaglar \href{https://ocw.mit.edu/courses/electrical-engineering-and-computer-science/6-254-game-theory-with-engineering-applications-spring-2010/index.htm}{[MIT 6.254]}
     \item Vincent Conitzer \href{https://courses.cs.duke.edu/spring16/compsci590.4/}{[Duke CPS 590.4]}
    \end{itemize}
    \item Elly Khodaie helped with importing slides from PowerPoint to \LaTeX
   \end{itemize}
  \end{frame}



\end{document}
